\documentclass{article}
\usepackage{amsmath}
\newcommand{\argmin}{\arg\!\min} 
\newcommand{\argmax}{\arg\!\max} 
\usepackage{amssymb}
\usepackage{amsfonts}
\usepackage[T1]{fontenc}
\usepackage{bm}
\usepackage{array}
\usepackage{graphicx}
\usepackage[utf8]{inputenc}

\title{Cart pulling research notes}
\date{\today}
\author{Marko Guberina}

\begin{document}
\pagenumbering{gobble}
\maketitle
\tableofcontents
\pagenumbering{arabic}

\section{Problem description}
\label{sec-problem-description}
We have a mobile robot with hands at our disposal. 
We want the robot to take a cart of things, and pull it from point A to B.
This is taking place in a crowded environment, so the robot needs to be able
to detect dynamic obstacles (people) and navigate around them.
The robot has lidar, a camera, positional encoders on everything, an IMU
and TODO: [other things we use].

The key aspects of this problem are:
\begin{enumerate}
		\item mapping and localization, as for anything mobile
		\item detecting obstacles
		\item very quick planning - we can't know where the people will go,
				so we need to plan quickly (can't rely on old plan which
				assumes obstacle follows a particular trajectory)
		\item controller capable of simultaneously following the 
				prescribed path and handling of the cart
\end{enumerate} 
We tackle each of these problem as follows:
\begin{enumerate}
\item Use off-the-shelf ROS2 Nav2 stack which has an implementation
	of SLAM. We know the environment we're in, and rely on a map made in 
	advance.
\item Some part of the ROS2 stack on the robot TODO: [know which one]
		takes the lidar point point cloud, filters out the unmapped 
		detected things as a separate layer on the map. Those can be assumed to be dynamic obstacles.
\item Use Albin Dahlin's [TODO: citation] dynamic system planning in star-world.
	Furthermore, use his code implementation. During operation,
	the planner's own map needs to be updated, as it relies on star-shaped
	obstacles, and not on a pixel-grid. This should use 
	as much caching as possible, but at least the static part of the map.
\item The ideal solution seems to be creating a whole-body MPC. 
Optimal control strategies can accommodate 
multiple optimization objectives, and underactuation in the base of the robot.
An alternative would be to disjointly control the base and the arms.
As will be explained later, we construct an exact path for the end-effector
to follow, so as long as it follows it, it doesn't matter whether
it's full or disjoint body control (as long as we don't have additional criteria
to on top of the base and the end-effector positions).
\end{enumerate}

\section{Solution approach}
\label{sec-solution-approach}

\subsection{Formal task description}
\subsection{System architecture}
TODO: finish

The only communication between the planner and the controller is 
the planner sending the path. 
However, in the current implementation, 
it was easier for the controller to communicate 
the current position to the planner.

\section{Control design}
\subsection{Optimal control formulation}
We propose a direct trajectory optimization approach for our conception
of the cart pulling problem. Furthermore, we solve the problem
on the dynamics level. We use the Box-FDDP solver [reference]
from the Crocoddyl optimal control suite [reference].
This approach has the following benefits:
\begin{itemize}
\item control of underactuated mobile bases like differential drive
\item flexibility in objective function design 
makes the approach extensible, for example allowing to penalization of centrifugal 
forces when handling tall carts, or extending the approach to legged robots
\end{itemize}

A generic optimal control problem looks like
%\begin{align*}
%		\label{eq-original-problem}
%\min_{x (\cdot), u (\cdot), T} &=
%\int_{0}^{T} L (x (t), u (t)) dt  + E (x (T))\\
%\text{subject to } &
%x (0) = x_{ 0 } \\
%& \dot{x} (t) = f (x (t), u (t)) \text{ (dynamics)}\\
%& h (x (t), u (t)) \geq 0 \text{ (path constraints)}\\
%& r (x (T)) =0 
%\end{align*}

\begin{align}
X^{ * }, U^{ * } &= \argmin_{X,U} \sum_{k=1}^{N-1} \int_{t_{ k }}^{t_{ k+1 }} l (x,u) dt +L (x (N)) \\
\text{such that } & \dot{x} = f (x,u)\\
& x \in \mathcal{X}, u \in \mathcal{U}\\
\end{align}

A common choice for $ \Delta t = t_{ k+1 } - t_{ k }  $ is
$ \Delta t = 0.01  $.


%We solve the inverse dynamics problem to solve for accelerations
%$ \tau = M(q)\ddot{q}+C(q,\dot{q})+G(q)$.
%TODO: what's the reason i can you either again????
%TODO: what's the reason i can you either again????
%TODO: what's the reason i can you either again????
%TODO: what's the reason i can you either again????
%TODO: what's the reason i can you either again????
% TODO: try both ,see if there's a difference.
% inv dyn should be better because no M(q) inversion, but idk

\subsubsection{Turning OCPs into MPCs}
Any OCP described here is turned into an MPC by:
\begin{itemize}
\item \textbf{Shortening the time horizon $ N  $, capping the maximum number of
solver iterations $ n_{ \text{solver\_iter} }  $, 
and reducing control frequency $ f_{ \text{ctrl} }  $}
until satisfactory performance is reached. 
From limited qualitative testing, the most common numbers
are $ N = 30  $, $  n_{ \text{solver\_iter} } = 10   $,
and $ f_{ \text{ctrl} } = 200  $.
\item \textbf{Warmstarting} with the previous solution. 
		The previous solution is offset by the amount of time passed.
		TODO: make sure it is, make sure the ctrl\_freq amount of time, 
		not $ \Delta t  $ time. If the second is the case,
		you probably want to linearly interpolate to offset correctly
		- also make sure you're time-testing whether this is actually helpful
		(if it doesn't shave off an iteration off of the solver,
		it isn't)
\end{itemize}

\subsubsection{Common OCP parts}
Depending on the robot used, and the task at hand,
there are different OCPs being used.
For example, to initialize the robot's position, we use an objective
function for point-to-point motions,
while we use a different one for path tracking.

Due to the complexity of the final controller, and of the total
control system design (including planning), during the development
process if was necessary to build smaller controllers
to increase understanding of the controller, informing 
and guiding the final design. 

It's easiest to begin by describing the common parts
of the all OCPs.
\subsubsection{Decision variables, dynamics contraint}
In our case, the state is a vector of joint positions and velocities:
\begin{equation}
 x = \begin{bmatrix} q & v \end{bmatrix}^{ T }   .
\end{equation}
The control inputs can be either acceleration or torques, 
depending on which we formulate either the inverse or the forward dynamics
problem respectively, and solve for the other quantity.
We used inverse dynamics.
In any case, because the states are also decision variables,
a state-input pair must follow the system dynamics into the next that.
Thus to put in the discrete version of $ \dot{x}= f (x,u)    $, we can write
\begin{gather}
 \text{dynamics} ( \delta x_{ k+1 }, \delta  x_{ k }, \delta  u_{ k } ) =0 \\
 \text{dynamics} (\delta \dot{x}, \delta x, \delta u)=
\delta \dot{x} - (f_{ x } \delta x + f_{ u } \delta u)
\end{gather}
TODO: elaborate further . Also explain to self whether efforts are torques.
also TODO: probably a good idea to actually formulate with the forward 
dynamics version too, it's about 10 basic lines of code and an argument.

\subsubsection{State constraints}
Since we use Box-FDDP, we can only have box constraints on torque inputs,
i.e. $ \tau_{ \min } \leq \tau \leq \tau_{ \max } $.
In the case of an underactuated joint, we set the torque limit on it to zero.

To put constraints on the state, 
namely on arms' joint limits and on all joints' velocities, 
we put them into a quadratic barrier function 
which we put into the objective function.
In particular (TODO make sure this is correct):
\begin{gather}
		r_{ lb } = x - x_{ \min } \\
		r_{ ub } = x - x_{ \max } \\
		B(x) =
\left\{ 
\begin{array}{ll}
		\frac{1}{2}  (\left\lVert r_{ lb } \right \rVert ^{ 2 }
		+ \left\lVert r_{ ub } \right \rVert ^{ 2 }) & \text{ if } 
		lb \leq x \leq ub \\
		\max \text{ float} & \text{ otherwise}
\end{array}
\right . 
\end{gather}

\subsubsection{Regulation costs (regularization)}
Finally, we add regulation costs on both the state $ c_{ 1 } x^{ T }x  $,
and control inputs $ c_{ 2 } u^{ T }u  $.

\subsection{Task-defining objective functions}
Depending on the task, we create a different objective function.
Before going into the functions themselves, the reader 
should be briefed about some specifics of Crocoddyl, the library
used for optimal control (more detailed text can be found in the appendix
of ``Robotics notes'').

At every timestep of the OCP, also known as a knot,
there is an ActionModel. An ActionModel store all data and functions
about the OCP at this point. This means that an OCP can be constructed
as a combination of different dynamics models, constraints, 
and objective functions, all defined just at one knot.
This is used in path-following, where every knot is assigned
a each own target (in this case a path point) to reach in the objective function.
Another useful feature is setting different cost coefficients,
which most useful in setting zero terminal velocities at the final knot.
Note that the last point is not a good idea in MPC!

\subsubsection{Point-to-point tasks}
In case there is a fixed end-goal to reach, we use the $se(3)$ error
norm as the task-specific function to put into the objective function.
\paragraph{End-effector pose task}
In the particular case of the goal being defined for the end-effector,
whose pose is denoted as $   T_{ w }^{ e }$ reaching a goal
denoted as $   T_{ w }^{ \text{goal} }$:
\begin{equation}
l_{ e } (x)= \left\lVert 
\log \left( \left( T_{ w }^{ e } \right)^{ -1 } \right) T_{ w }^{ \text{goal} }
\right \rVert _{ 2 }
\end{equation}
Here we of course use only joint values $ q  $, not the full state $ x  $,
because the end-effector pose is only a function of $ q  $.

We put this cost function at every knot.

\paragraph{End-effector pose and base position}
In case we also want to specify the position of the base,
we additionally put in the translational cost for the base.
If we denote the pose of the base is 
\begin{equation}
T_{ w }^{ b } = \begin{bmatrix} 
		R_{ w }^{ b } & p_{ b } \\
		0 & 1
\end{bmatrix} 
\end{equation}
then the translational cost is given by
\begin{equation}
		l_{ b } (x) = \left\lVert   p_{ \text{goal} } - p_{ b } \right \rVert
\end{equation}
Here it should be noted that the $ z  $ coordinate is 0 for
both the goal and the base pose.
The orientation of the base could be accounted for in the same manner.

To get the joined task, we simply sum the two goal costs:
\begin{equation}
l_{ \text{joined} } (x)= l_{e } (x) + l_{ b } (x)
\end{equation}

\subsubsection{Path-following tasks}
As discussed in \ref{sec-problem-description}, the controller is tasked with following 
a path prescribed by a path planner.
One way, and probably the best way,
to do this would be to parametrize the path with a parameter $ s \in [0,1] $, 
and task the controller
with timing it, i.e. constructing a trajectory on it. In this case the 
time-scaling factor of the path $ \lambda  $ in the equation
$ t = \lambda s   $ would be a decision variable in the OCP.

However, we found it easier to manually construct a trajectory, i.e. a timing law
on the path. This is because it is very easy to construct a trajectory following OCP.
Namely, it is done by modifying the point-to-point formulation by replacing
$ T_{ w }^{ goal }  $ with $ T_{ w }^{ traj } [i]  $ where 
$   T_{ w }^{ traj } [i]    $ is the point on the desired trajectory corresponding
to the desired timing law on the path. More on path and trajectory construction can
be found in \ref{sub-reference-traj-construction}.
The important point here is that the trajectory needs
to be interpolated so as to correspond to the timing of the knots.
To be more specific, the time difference $ \Delta t  $
between two indexes of the array defining the reference trajectory
has to be the one in the OCP.
For our purposes linear interpolation is sufficient,
especially because the time-horizon in the MPC is rather short,
leading to a short trajectory in space as well

\paragraph{End-effector trajectory following cost function}
The task-defining running cost for end-effector trajectory is 
given by (here with explicit indexing, ):
\begin{equation}
l_{ path } (x [i]) = 
\left\lVert \log 
\left( \left( T_{ w }^{ e } [i] \right) _{ 1 } \right) T_{ w }^{ \text{ref} } [i]
\right \rVert_{ 2 } 
\end{equation}

\paragraph{Base path following cost function}
Works in the same way as the described above for the end-effector case,
just with different translations instead.

\paragraph{Base and end-effector trajectory following cost function}
Is again simply the sum of the two costs.

\subsubsection{Solver choice}
Box-FDDP from Crocoddy seems to be the fastest, 
and provides control input constraints
which is crucial - otherwise one would need to overtune control input regulation,
leading to a steep decrease in performance.

CSQP from mim-solvers was also tested, but it was too slow (NOTE: idk
if I set it up right, the paper said it should be faster than what I got).
The benefit it has is that it can accommodate other linear constraints
(it's based on sequential quadratic programming (SQP), made faster
for trajectory optimization through it's implementation).

Casadi' IPOPT is a good choice for experimenting 
with different problem formulations,
but being a generic optimization solver, and all the derivatives 
being obtained through finite differencing, it is of course too slow 
for practical implementation.


\subsection{Reference trajectory construction}
\label{sub-reference-traj-construction}
\subsubsection{Base reference trajectory}
The path planner produces a path 
$ \gamma_{ 2 }  \in \mathbb{R}^{2} $ 
for the base of the robot.
It is given by an array of points $ \Gamma \in \mathbb{R}^{n \times 2} [i] $
where $ i  $ is the index denoting $ i  $th element (point) of the array,
and indexing starts with 0.
We begin by introducing a physically dimensionless parametrization 
$ s \in [0,1]  $, resulting in $ \gamma_{ 2 } (s)  $.
Importantly, we assume that $ \Gamma  $ is constructed
from a uniform sampling of $ \gamma  $ (this is true with the
chosen path planner).

We want to impose a timing law so that the path is traversed at a fixed
velocity, $ v_{ \text{ref} }  c_{ 3 } \max v_{ \text{base} }  $ 
where $ c_{ 3 } \in [0,1]  $.
A different timing law ensuring the same velocity will be constructed
for the end-effector reference.

We construct the timing law by calculating the time required
to traverse the prescribed path at this velocity.
To do so we first calculate the arclength of the path via
$ L (\gamma) = \int_{s}^{} \left\lVert    \dot{\gamma} (s)\right \rVert  ds  $.
In discrete space perform the following computation
$ L (\Gamma) =  \sum_{i=1}^{n} \left\lVert 
\Gamma [i+1] - \Gamma [i]  \right\rVert $.
From this we get 
$ t_{ \text{final} } = \frac{  L (\gamma) }{ v_{ \text{ref} }}  $,
leading to the timing law $ t = t_{ \text{final} } s  $,
and the reference trajectory $ \gamma (t (s))  $.

\subsubsection{End-effector reference trajectory}
As discussed in \ref{sec-solution-approach}, we want
the cart to follow the past path of the base.
We denote the continuous past path with 
$ \phi (s) \in \mathbb{R}^{2}, s \in [0,1]  $, and 
$ \Phi [i] \in \mathbb{R}^{n \times 2}  $ as the discrete collection of points.
In the software implementation, $ \Phi  $ is a double ended queue
where at every iteration of the control loop a new
point $ p_{ \text{base} }  $ is inserted, and the oldest
entry is removed. We refer to this double ended queue as the rolling past window.
The size of the rolling past window $ n  $ is set by the user,
and it should be long enough to construct the preferred trajectory.
Note that $ \Phi  $ is initialized as a straight line from the end-effector
to the base of the robot. 

Assuming a stable control frequency (it empirically holds well enough),
$ \Phi  $ is a uniform sampling of $ \phi  $.
We want to maintain a fixed distance between the base and the handlebar of
the cart. Since the handlebar is a rigid body, it also uniquely
determines the references for the end-effector(s).
Thus the first task is to find the point on the past path
which is the preferred distance from the robot on the arclength 
of the past path. As in the case of calculating the length
of the path given by the path planner for the base, 
we use 
$ L (\Phi) = \sum_{i=1}^{n} \left\lVert 
\Phi [i+1] - \Phi [i]
\right \rVert   $. 
By computing this sum in a loop, we find the index $ h  $
for which $   $
$ L (\Phi) = \sum_{i=1}^{h} \left\lVert 
\Phi [i+1] - \Phi [i]
\right \rVert  = d_{ \text{preferred} }  $,
where $ d_{ \text{preferred} }  $ denotes the preferred
distance of the handlebar to the base of the robot.
Since we do not care about the past path beyond the index $ h  $,
we cut them off of $ \Phi  $. The resulting
path from the preferred handlebar distance to the bast of the robot
can then be parametrized with $ s \in [0,1]  $.

The next step is to construct an $ SE (3)  $ reference from the 2-dimensional
past path in $ (x,y)  $ coordinates of the world frame.
In the dual arm case this will be the absolute frame
of the dual arm manipulation task.
Since the robot is operating on flat ground, and the cart is rigid,
the height, i.e. the $ z  $ coordinate is known in advance and fixed.
Since the handlebar is rigidly grasped, and we want the cart
to follow the base's path, the rotational axes are fixed as well.
We set them by constructing a Frenet frame on the past path.
Without loss of generality, let the $ z  $ axis of the end-effector frame
be perpendicular to the ground, and pointing down.
Also let the $ x  $ axis of the end-effector frame be parallel
to the tangent of the path (either direction is fine,
we choose toward the mobile base). This ensures the same orientation
of the cart as the base of the robot on the path.
Finally, the $ y   $ axis simply completes the 
right-handed frame.

To make this construction, we need to compute the tangent of the past path 
$ \dot{\phi} (s)  $. Alternatively base velocities could also be logged,
but this approach is less prone to noise.
TODO: this is not a 100\% obvious statement due to potential
localization-induced teleporting.
For simplicity, we compute the angle of the secant of the path
in the world frame, and formulate the rotation matrix from it.
In particular 
\begin{equation}
 \theta [i] = atan2 \left(  \frac{\phi_{ y } [i+1] - \phi _{ y } [i]}{\phi_{ x } [i+1] - \phi _{ x } [i]}  \right)  
\end{equation}
 which leads to
 \begin{equation}
 R_{ w }^{ \text{ref} } [i] = 
\begin{bmatrix} 
		\cos (\theta [i]]) &  \sin (\theta [i]) & 0  \\
		\sin (\theta [i]]) & - \cos(\theta [i]) & 0  \\
		0 & 0 & -1  
\end{bmatrix} 
 \end{equation}









\end{document}
